\chapter{RISC-V}\label{chap:riscv}

\section{Overview}

RISC-V is an open source instruction set architecture first published in May 2011 by A. Waterman et al. \cite{first_riscv}.
As contained in the name, it is based on the RISC design philosophy. \todo{explain risc (compare wiki)}
Since 2015 the developement of RISC-V is coordinated by the RISC-V International Association, a non-profit corporation based in switzerland since 2020 \cite{riscvorg}.
Its goals are among others an \emph{open} ISA that is freely available to all,
a \emph{real} ISA suitable for native hardware implementation and
an ISA separated into a \emph{small} base integer ISA usable by itself e.g. for educational purpose
and optional standard extensions to support general purpose software development \cite{riscv-isa}(Chapter 1).

It currently contains four base ISAs, namely RV32I, RV64I, RV32E and RV64E,
wich may be extended with one or more of the 47 ratified extension ISAs \cite{riscv-isa} (Preface).

\extend{Hier brauchts vlt noch was} \todo{little endian erwähnen?}

For my work, I will focus on a subset of the RV64I ISA.

\section{The RV64I ISA}
RV64I is not complex, but its structure is relevant to understand my later work.
So I explain all elements relevant for my thesis.

RV64I has 32 64bit registers called $x0$-$x31$, where $x0$ is hardwired to 0 on all bits.
Registers $x1$-$x31$ are general purpose and are interpreted by various instructions as a collection of booleans, two's complement signed binary integers or unsigned integers.
Additionaly there is a register called $pc$ acting as program counter and holding the address of the current instruction \cite{riscv-isa}(Chapters 4.1, 2.1).

In RV64I, a memory address has the size of 64bit. As the memory model is defined to be single byte adressable,
the address space of RV64I is $2^{64}$ bytes \cite{riscv-isa}(Chapter 1.4).

Like almost all of the standard ISAs of RISC-V, RV64I has a standart encoding lenght of 32bit or 1 \emph{word}.
Only the compressed extension C adds instuctions with a length of 16bit \cite{riscv-isa}(Chapter 1.5), but it is irrelevant for us.
All instructions of RV64I are encoded in one of the six formats shown in \figref{fig:rv64i_formats}.\\
The design of these format results in the following features:
\begin{itemize}
    \item As of RISC-Vs little endianess, the $opcode$, wich encodes the general instruction, is always read first.
            Also, further specification of the instruction by $funct3$ and $funct7$ is found always at the same location.
    \item If used by the instruction, the destination register $rd$ and the source registers $rs1$ \& $rs2$ are always at the same place. This simplifys decoding.
    \item The highest bit of the immediate value $imm$ is always at the last bit. This makes finding the sign of a signed immediate value trivial.
\end{itemize}

\begin{figure}[t]
    \begin{centering}
        \includegraphics[width = \textwidth]{figures/2-RiscV/R.png}\\
        \includegraphics[width = \textwidth]{figures/2-RiscV/I.png}\\
        \includegraphics[width = \textwidth]{figures/2-RiscV/S.png}\\
        \includegraphics[width = \textwidth]{figures/2-RiscV/B.png}\\
        \includegraphics[width = \textwidth]{figures/2-RiscV/U.png}\\
        \includegraphics[width = \textwidth]{figures/2-RiscV/J.png}
        \caption[RV64I encoding formats]{RV64I encoding formats, used in \cite[Chapter 2.3]{riscv-isa}\todo{Kopie richtig angeben}}
        \label{fig:rv64i_formats}
    \end{centering}
\end{figure}





\section{Simulation of RISC-V}

\subsection{Saving the State of a RISC-V Processor}\label{statefile}